% Section 1 -- introduction, half a page. Selection framing, three
% contributions, one statement of the range. The abstract and 5.1 carry the
% rest. No literal numbers: macros only.

\section{Introduction}
\label{sec:intro}

A Lean proof the kernel accepts is correct, and nothing more needs saying about
that. But a proof is also a \emph{selection}. Of the hundreds of thousands of
declarations in a library, it reaches for a particular few. That selection
determines what a premise-selection model must learn to retrieve, what breaks
when a lemma is renamed or generalized, and whether a development participates
in the library or merely sits beside it. Correctness is settled by the kernel. Nothing settles the selection.

We make it measurable. A proof's premises are nodes in the library's dependency
graph, so a proof population induces a subgraph, and two populations can be
compared by which parts of the library their subgraphs occupy. We compare a
released machine corpus of \machineProofs{} Lean proofs against mathlib's
\humanProofs{} declarations, working entirely from published artifacts and
invoking Lean nowhere.

Three contributions follow. First, what a released-artifact pipeline can
measure and how reliably (Section~\ref{sec:instrument}): the release's
\texttt{is\_explicit} flag does not mean ``written in source'', so no
source-level instrument should be expected to reach unit retention against it;
we report retention as a scope decomposition at \retention{}, of which
\efficiencyReachable{} is what a per-declaration instrument can reach, flat
across tactic classes so the instrument does not fail preferentially on the
proof styles that separate the corpora; we quantify the displacement our own
extractor induces, and show that proof length must be measured on a
wrap-invariant axis. Second, the coverage result
(Section~\ref{sec:results}): mathlib's proofs draw on \humanVocab{} distinct
premises against the machine corpus's \machineVocab{}, a gap that survives
every control we can construct at \vocabRatio{}$\times$ on the full corpora,
\areaRestrictedRatio{}$\times$ restricted to the machine vocabulary's own
library regions, and \vocabRatioMatched{}$\times$ matched on proof count,
tactic class and length. The direction is invariant across designs and the
magnitude is not, so we quote the range throughout. Third, a caution
(Section~\ref{sec:dag-layer}): a depth finding we built on a published
attribute is withdrawn, because the attribute is a centrality proxy, and we
give the check that would have caught it.
